\section{Items}\label{sec:items}\index{Item}\index{Weapon|see{Item}}\index{Armor|see{Item}}\index{Shield|see{Item}}
Beyond a character's own abilities, they are also generally equipped with certain equipment that aids in their adventures.
Much of the equipment is not noteworthy and can generally be assumed, but the key pieces of equipment are \textbf{Items} and given their own statistics.
These statistics define bonuses and other features that the \emph{Item} may grant to the character using it.
\emph{Items} fall into one of two categories: \textbf{Weapons} and \textbf{Relics}.
Weapons, including defensive items like shields and armor, are associated with a particular Combat Style the character has and grant bonuses to combat abilities.
Relics are any non-Weapon items and they grant the user an additional \Experience{} that reflects the nature or possibly history of the Relic.
Both Weapons and Relics can also grant other benefits, including bonuses to other \emph{Attributes} or \emph{Specializations}, additional \emph{Experiences}, or even new \emph{Abilities}.

\subsection{Creating an Item}
At character creation a character receives one Item for free, but otherwise creating an Item costs 1 XP which reflects the costs of time, money, and resources that procuring or creating the Item takes.
For this cost, the item that is created is a \emph{Basic Item}, the properties of which depend on whether it is a \emph{Weapon} or \emph{Relic}.
When the item is created you can also spend XP to add additional features or improve on the Items basic features.

\subsubsection{Creating a Basic Weapon}\index{Item!Weapon}
To create a basic weapon, choose one of the character's Combat Styles to associate it with.
A basic weapon grants a $+1$ bonus to attack and damage rolls when used with an Offensive Combat Style and a $+1$ to defend rolls and a $+1$ bonus to \Block when used with a Defensive Combat Style.

\subsubsection{Creating a Basic Relic}\index{Item!Relic}
To create a basic relic, define the \Experience{} that the Relic grants a $+1$ bonus for.
At the beginning of each session, the Relic gains an \Advantage{} (if it does not have one already) that can be spent as usual.

\subsubsection{Finishing an Item}
Finally, additional features can be purchased for an Item using XP.
\begin{itemize}
  \item A bonus to an \Attribute{}, to \Block{}, to attack/defend (only for weapons), or a new \Experience{} (or to increase the bonus to a relic's \Experience{}).
    The cost of this bonus is based on the \emph{total bonus} the item offers (this is the total bonus for attacking, defending, to \emph{Attributes}, to \emph{Block}, to \emph{Experiences}, and to \emph{Specializations}).
    For basic statistics and Block (which normally cost $3 \times$ their new bonus), the cost is the total bonus plus $2\times$ the total bonus to that stat or the total bonus to Block, respectively.
    For \emph{Experiences} and bonuses to attack/defend the cost is just the new total bonus, and this is also the cost for adding a new $+1$ \Experience.
    As an example, if we have a basic weapon, which grants $+1$ to defend and $1$ block, increasing the bonus to defend to $+2$ will cost $3$ XP, and then adding a $+1$ bonus to \emph{Strength} costs 6 XP (4 because the total bonus is now $+4$ and 2 more because the total bonus to Strength is $+1$ and the cost includes $2 \times$ that bonus).
    The maximum bonus to an \Attribute{}, \Block, \Experience{}, or to attack or defend is $+2$ at character creation and $+3$ thereafter.
  \item A \Specialization{} bonus can either offer just a bonus to one of the character's existing \emph{Specializations} or can offer a new \Specialization{}.
    The cost of a \Specialization{} bonus is the same as described above (i.e., the total bonus the item offers after the new bonus) and if it provides a new \emph{Specialization} then it also costs whatever tag it provides.
  \item A bonus to \Life{} costs 2 XP per point up to a maximum bonus of $+ 10$.
  \item A bonus to \Recovery{} costs 4 XP per point up to a maximum bonus of $+ 5$.
  \item Granting an \Ability{} costs the cost of that \Ability{}.
\end{itemize}
For weapons, the following options can also be added
\begin{itemize}
  \item An additional Combat Style that the weapon can be used with costs 3 XP.
  \item The weapon can add additional damage to attacks using it, with the number of XP depending on the damage die added: 5 XP for a d4, 15 XP for a d6, 30 XP for a d8, 50 XP for a d10, and 75 XP for a d12 (an existing die can be increased to the next higher die size for the difference in their XP costs).
\end{itemize}

\subsection{Other Item Rules}\index{Item}
There are a few other Item rules, especially for governing how they interact.
First, if a character has two Items which grant a bonus to the same thing (the same \Attribute{} or \Specialization) these bonuses do not stack, instead only the higher bonus applies.
Similarly, if a character has multiple different weapons, the attack bonus only applies when using that specific weapon.
For defensive weapons, similarly only the highest bonus to defend applies.

Additionally, if a character has more than 2 items in use (a character can hold onto other item but if they are not in use they would require several hours to be able to use them) they must hold onto 1 XP for each item past 2 they have.

Finally, sometimes a character outgrows a particular item, either they find something better or they have changed enough that the item no longer makes sense for them.
In such case they can sell the item to regain all XP that they had invested in creating the item (make sure to track this for all Items); because this represents, in game, finding a buyer or perhaps melting down a weapon to forge it into something else.
The GM may not allow an Item to be sold in a circumstance where it does not fit the narrative.
